%%%%%%%%%%%%%%%%%%%%%%%%%%%%%%%%%%%%%%%%%%%%%%%%%%%%%%%%%%%%%%%%%%%%%%%%%%%%%%
\section{The Physical Problem}
%%%%%%%%%%%%%%%%%%%%%%%%%%%%%%%%%%%%%%%%%%%%%%%%%%%%%%%%%%%%%%%%%%%%%%%%%%%%%%

Relativistic quantum mechanics and quantum field theory provide an extremely
successful description of elementary particles. Nevertheless, from the
viewpoint of dynamics they differ fundamentally from classical relativistic
mechanics. In the latter, the state of a particle at any instant is specified
by its position and momentum, and the subsequent evolution follows from the
equations of motion. In quantum theory, by contrast, the state is represented
by a wave function or a state vector, while the notion of an individual
trajectory no longer occupies a fundamental role.

An obvious question therefore arises. Is it possible that the quantum
description is itself a statistical description of a more fundamental
relativistic dynamics? More specifically, can the motion of a relativistic
particle be described by a stochastic process whose statistical properties
reproduce the equations of relativistic quantum mechanics?

In the nonrelativistic domain this question was answered, at least partially,
by the stochastic mechanics of Fényes and Nelson. Starting from a diffusion
process in ordinary three-dimensional space, Nelson showed that under suitable
assumptions the probability density and current satisfy equations equivalent
to the Schrödinger equation. Whether one accepts this interpretation or not,
it demonstrates that quantum dynamics can arise from an appropriately chosen
stochastic process.

The relativistic case is considerably more difficult. The difficulty is not
primarily the construction of a relativistic diffusion process. Various forms
of relativistic Brownian motion have been proposed during the past sixty
years. Rather, the central problem is to construct a stochastic dynamics that
simultaneously satisfies the following physical requirements:

\begin{enumerate}

\item
Manifest Lorentz covariance.

\item
A well-defined Hamiltonian dynamics.

\item
Positive probability densities.

\item
Compatibility with gauge interactions.

\item
A correct nonrelativistic limit.

\item
Reduction to the experimentally verified equations of relativistic quantum
mechanics.

\end{enumerate}

No presently known theory satisfies all of these requirements simultaneously.
Some approaches possess manifest covariance but fail to reproduce quantum
mechanics. Others reproduce the Schrödinger equation but rely on a preferred
time coordinate and therefore lose manifest covariance. Still others introduce
relativistic diffusion but abandon the Hamiltonian structure that underlies
both classical mechanics and the SHP formalism.

The purpose of this monograph is to investigate whether these apparently
independent difficulties may admit a common resolution within the framework of
manifestly covariant dynamics.

To understand why this possibility deserves serious consideration, it is
necessary to examine where the standard relativistic generalization of
Brownian motion first encounters difficulties.
%%%%%%%%%%%%%%%%%%%%%%%%%%%%%%%%%%%%%%%%%%%%%%%%%%%%%%%%%%%%%%%%%%%%%%%%%%%%%%
\section{The First Obstacle: Relativistic Brownian Motion}
%%%%%%%%%%%%%%%%%%%%%%%%%%%%%%%%%%%%%%%%%%%%%%%%%%%%%%%%%%%%%%%%%%%%%%%%%%%%%%

Consider first ordinary Brownian motion in Euclidean space. During a small
time interval $dt$, the random displacement satisfies
\[\langle dW_i \rangle =0,\]
and
\[\langle dW_i\,dW_j\rangle=2D\,\delta_{ij}\,dt,\]
where $D$ is the diffusion coefficient.
The covariance matrix is positive definite, guaranteeing that the variance of
every component is non-negative.

The obvious relativistic generalization would appear to be
\[\langle dW^\mu\rangle =0,\]
and
\[\langle dW^\mu dW^\nu\rangle=2D\,\eta^{\mu\nu}\,d\tau,\]
where $\eta^{\mu\nu}$ is the Minkowski metric.
This expression is manifestly Lorentz covariant. Unfortunately, it cannot
describe an ordinary stochastic process.
Indeed,
\[\langle (dW^0)^2\rangle=2D\,d\tau,\]
whereas
\[\langle (dW^1)^2\rangle=-2D\,d\tau,\]
and similarly for the remaining spatial components. The variances of the
spatial increments are therefore negative, which is impossible for real random
variables.

This elementary observation constitutes the first and perhaps most fundamental
difficulty of relativistic stochastic mechanics. The problem appears before
any dynamical equation has been written and is independent of the particular
choice of interaction or Hamiltonian. It follows directly from the indefinite
signature of Minkowski space.

Consequently, relativistic stochastic mechanics cannot be obtained by the
simple replacement
\[\delta_{ij}\longrightarrow\eta_{\mu\nu}.\]
A successful relativistic theory must therefore modify the very notion of
diffusion rather than merely rewriting the nonrelativistic equations in a
covariant notation.

%%%%%%%%%%%%%%%%%%%%%%%%%%%%%%%%%%%%%%%%%%%%%%%%%%%%%%%%%%%%%%%%%%%%%%%%%%%%%%
\section{The Second Obstacle: The Evolution Parameter}
%%%%%%%%%%%%%%%%%%%%%%%%%%%%%%%%%%%%%%%%%%%%%%%%%%%%%%%%%%%%%%%%%%%%%%%%%%%%%%

The failure of the naive relativistic generalization of Brownian motion
discussed in the previous section demonstrates that relativistic stochastic
mechanics cannot be obtained by simply replacing Euclidean geometry with
Minkowski geometry. Before attempting to construct a relativistic diffusion
process, one must first answer a more elementary question:
\begin{center}
\emph{With respect to which parameter should the stochastic evolution be
defined?}
\end{center}

In nonrelativistic mechanics the answer is obvious. The trajectory of a
particle,
\[{\bf x}={\bf x}(t),\]
is parametrized by the universal Newtonian time $t$. Since every observer
agrees on the value of $t$, it provides a natural ordering of successive
events. Consequently, the stochastic process
\[{\bf X}(t)\]
is defined without ambiguity.

Relativistic mechanics changes this situation completely. The coordinate time
\[x^0=ct\]
is no longer an invariant quantity. Different inertial observers assign
different time coordinates to the same event, related by Lorentz
transformations,
\[x'^{0}=\gamma(x^{0}-\beta x^{1}).\]
Thus a stochastic process defined with respect to $x^0$ is tied to a
particular Lorentz frame and therefore loses manifest covariance.

One might attempt to replace coordinate time by the proper time
\[ds^{2}=dx^{\mu}dx_{\mu}.\]
Although proper time is Lorentz invariant, it is not suitable as a universal
evolution parameter. It is defined only along timelike worldlines and depends
upon the trajectory itself. Furthermore, interacting systems containing
several particles possess, in general, several different proper times.
Consequently, proper time cannot provide a common dynamical parameter for a
many-body relativistic theory.

These observations indicate that neither coordinate time nor proper time
possesses all the properties required for a stochastic evolution parameter.
The SHP formalism resolves this difficulty by introducing an invariant
evolution parameter
\[\tau,\]
which is independent of the space-time coordinates. The particle trajectory is
therefore written as
\[x^{\mu}=x^{\mu}(\tau),\]
where all four coordinates are treated as dynamical variables.

Unlike coordinate time, $\tau$ is invariant under Lorentz transformations.
Unlike proper time, it is independent of the particular trajectory and may be
used simultaneously for interacting many-particle systems. It therefore
provides a universal ordering parameter for relativistic dynamics.

From the viewpoint adopted in the present monograph, this observation is of
fundamental importance. The principal achievement of the SHP formalism is not
merely the introduction of an additional parameter, but the identification of
a physically meaningful invariant quantity possessing exactly those properties
required for the formulation of a relativistic dynamical theory.

If one wishes to construct a stochastic extension of relativistic mechanics,
it is therefore natural to replace the deterministic trajectory
\[x^{\mu}=x^{\mu}(\tau)\]
by a stochastic trajectory
\[X^{\mu}=X^{\mu}(\tau,\omega),\]
where $\omega$ labels individual realizations of the stochastic process.
The remainder of this monograph is devoted to investigating whether such a
construction can reproduce the known structure of relativistic quantum
mechanics.
%%%%%%%%%%%%%%%%%%%%%%%%%%%%%%%%%%%%%%%%%%%%%%%%%%%%%%%%%%%%%%%%%%%%%%%%%%%%%%
\section{Why Manifestly Covariant Hamiltonian Dynamics?}
%%%%%%%%%%%%%%%%%%%%%%%%%%%%%%%%%%%%%%%%%%%%%%%%%%%%%%%%%%%%%%%%%%%%%%%%%%%%%%

The introduction of an invariant evolution parameter solves only one part of
the problem. A complete relativistic dynamics also requires equations of
motion. The simplest possibility would be to parameterize ordinary relativistic
mechanics by the invariant parameter $\tau$. Such a construction, however,
would merely constitute a reparametrization of the conventional theory and
would not provide an unconstrained Hamiltonian dynamics.

The difficulty originates from the standard relativistic relation
\[p^\mu p_\mu = m^2c^2,\]
which is usually imposed as a constraint. The Hamiltonian of the system
vanishes identically and the equations of motion are generated by the
constraint itself. This formulation is entirely adequate for many classical
problems, but it is not well suited for the construction of a stochastic
theory.

A stochastic process requires a genuine evolution equation with respect to an
independent parameter. If the evolution parameter itself is determined by a
constraint, the stochastic interpretation becomes obscure. It is therefore
desirable to formulate relativistic mechanics as an unconstrained Hamiltonian
system whose evolution is generated by an ordinary Hamiltonian function.

This is precisely the approach introduced independently by Stueckelberg and
subsequently developed by Horwitz and Piron.
The fundamental dynamical variables are
\[x^\mu(\tau), \qquad p_\mu(\tau),\]
while the evolution is generated by a Hamiltonian function
\[K(x,p).\]

The canonical equations assume exactly the same form as in ordinary
Hamiltonian mechanics,
\[\frac{dx^\mu}{d\tau}=\frac{\partial K}{\partial p_\mu},\]
\[\frac{dp_\mu}{d\tau}=-\frac{\partial K}{\partial x^\mu}.\]
The only formal difference is that the evolution parameter is now the
Lorentz-invariant quantity $\tau$ rather than Newtonian time.

This observation has far-reaching consequences. The mathematical structure of
relativistic dynamics becomes almost identical to that of ordinary Hamiltonian
mechanics. Canonical transformations, Poisson brackets, Liouville's theorem,
Hamilton-Jacobi theory, and statistical mechanics all carry over with only
minor modifications. From the standpoint of dynamics, manifest covariance is
achieved without abandoning the familiar Hamiltonian framework.

For a free particle the Hamiltonian is chosen as
\[K_0=\frac{p^\mu p_\mu}{2M},\]
where $M$ is a constant parameter with the dimensions of mass. Interactions
are introduced through an invariant potential,
\[K=\frac{p^\mu p_\mu}{2M}+V(x).\]
The corresponding equations of motion are
\[\frac{dx^\mu}{d\tau}=\frac{p^\mu}{M},\]
and
\[\frac{dp_\mu}{d\tau}=-\frac{\partial V}{\partial x^\mu}.\]

These equations are structurally indistinguishable from Newtonian Hamiltonian
mechanics. Consequently, every construction based upon Hamiltonian flows may,
at least in principle, be generalized to the relativistic domain.

This observation is central to the present monograph. Our objective is not to
construct a stochastic process directly in Minkowski space. Rather, we seek a
stochastic extension of an already well-defined relativistic Hamiltonian
dynamics.

From this perspective, the SHP formalism plays a role analogous to that of
Newtonian mechanics in Nelson's original construction. It supplies the
deterministic dynamics upon which stochastic fluctuations may subsequently be
superimposed.
%%%%%%%%%%%%%%%%%%%%%%%%%%%%%%%%%%%%%%%%%%%%%%%%%%%%%%%%%%%%%%%%%%%%%%%%%%%%%%
\section{From Deterministic to Stochastic Dynamics}
%%%%%%%%%%%%%%%%%%%%%%%%%%%%%%%%%%%%%%%%%%%%%%%%%%%%%%%%%%%%%%%%%%%%%%%%%%%%%%
Having established a manifestly covariant Hamiltonian dynamics, we may now
address the central question of this monograph.
How should stochastic effects be incorporated into the equations of motion?

It is important to appreciate that this question differs fundamentally from
the construction of Brownian motion in Euclidean space. The deterministic
dynamics already exists. The objective is therefore not to invent an entirely
new dynamics but rather to determine how random fluctuations modify an
existing Hamiltonian evolution.

This viewpoint is familiar throughout theoretical physics. Classical kinetic
theory begins with Newton's equations and subsequently introduces statistical
descriptions of large ensembles. Langevin dynamics begins with Newtonian
mechanics and supplements it by fluctuating and dissipative forces. Likewise,
the Fokker--Planck equation describes the statistical evolution associated
with an underlying stochastic modification of deterministic trajectories.

The same philosophy will be adopted here.
The deterministic motion generated by the SHP Hamiltonian will be regarded as
the zeroth-order approximation. Stochastic dynamics should then arise as a
systematic extension rather than as an independent postulate.
Accordingly, the deterministic equations
\[\frac{dx^\mu}{d\tau}=\frac{\partial K}{\partial p_\mu},\]
\[\frac{dp_\mu}{d\tau}=-\frac{\partial K}{\partial x^\mu},\]
will be replaced by stochastic evolution equations of the schematic form
\[dx^\mu=\frac{\partial K}{\partial p_\mu}\,d\tau+d\xi^\mu,\]
\[dp_\mu=-\frac{\partial K}{\partial x^\mu}\,d\tau+d\eta_\mu.\]

At the present stage no assumptions are made concerning the statistical
properties of the fluctuations
\[d\xi^\mu,\qquadd\eta_\mu.\]
In particular, they need not represent ordinary Brownian motion.
Their precise mathematical structure constitutes one of the principal
questions investigated in the subsequent chapters.

The advantage of this formulation is immediately apparent. The deterministic
Hamiltonian dynamics remains explicitly visible throughout the construction.
Randomness appears only through additional terms describing fluctuations about
the underlying Hamiltonian flow.

Consequently, every successful result obtained in deterministic SHP dynamics
remains available. Gauge interactions, many-particle systems, canonical
transformations and Hamilton-Jacobi theory require extension rather than
replacement.
%%%%%%%%%%%%%%%%%%%%%%%%%%%%%%%%%%%%%%%%%%%%%%%%%%%%%%%%%%%%%%%%%%%%%%%%%%%%%%
\section{The Origin of Stochasticity}
%%%%%%%%%%%%%%%%%%%%%%%%%%%%%%%%%%%%%%%%%%%%%%%%%%%%%%%%%%%%%%%%%%%%%%%%%%%%%%

The equations proposed in the preceding section introduce random
increments into the deterministic SHP dynamics. At this stage,
however, the origin of these fluctuations has not been specified.
This omission is intentional.

The mathematical construction of stochastic dynamics does not require
an explicit microscopic model. Indeed, in classical statistical
mechanics the Langevin equation provides an accurate description of
Brownian motion without referring to the detailed dynamics of every
molecule of the surrounding fluid.

Nevertheless, the physical interpretation of the stochastic terms
cannot be ignored indefinitely. A relativistic stochastic mechanics
should eventually explain why the motion of an isolated microscopic
particle differs from the deterministic evolution predicted by
classical Hamiltonian mechanics.

Several possibilities may be distinguished.

First, the stochastic forces may represent the influence of
unobserved microscopic degrees of freedom belonging to the physical
vacuum.

Second, they may arise from internal dynamical degrees of freedom of
the particle itself, which have been averaged over in the effective
description.

Third, they may represent an effective description of interactions
with additional fields not explicitly included in the Hamiltonian.

Finally, stochasticity may be regarded as a purely phenomenological
description whose justification lies entirely in its ability to
reproduce experimentally observed quantum phenomena.

The present monograph deliberately adopts the last viewpoint during
its initial development. No specific microscopic mechanism will be
assumed. Instead, the statistical properties of the stochastic
increments will be determined by consistency with relativistic
dynamics and with quantum mechanics.

Only after the mathematical structure has been established shall we
return to the question of the physical origin of the fluctuations.

